\documentclass[dvipdfmx,11pt]{jsarticle}

\usepackage{amsmath, amssymb, amsthm}
\usepackage{ascmac}
\usepackage{color}
\usepackage{empheq}
\usepackage{mathrsfs}
\usepackage[margin=25truemm]{geometry}
\usepackage{setspace}

\usepackage{tikz}
\usetikzlibrary{calc, angles, quotes}




\newcommand{\Div}{\operatorname{div}}
\newcommand{\dsp}{\displaystyle}
\newcommand{\e}{\varepsilon}
\newcommand{\Ha}{\mathcal{H}}
\newcommand{\prob}{\vspace{2mm}\pro}
\newcommand{\redprob}{\vspace{2mm}\color{red}\pro\color{black}}
\newcommand{\blueprob}{\vspace{2mm}\color{blue}\pro\color{black}}
\renewcommand{\qedsymbol}{$\blacksquare$}
\newcommand{\R}{\mathbb{R}} 
\newcommand{\setmid}{\mathrel{}\middle|\mathrel{}}


\theoremstyle{definition}
\newtheorem{theorem}{Theorem}[section]
\newtheorem{lemma}[theorem]{Lemma}
\newtheorem{corollary}[theorem]{Corollary}
\newtheorem{Def}[theorem]{Definition}
\newtheorem{prop}[theorem]{Proposition}
\newtheorem{rem}[theorem]{Remark}
\newtheorem{ex}[theorem]{Example}
\newtheorem{assumption}{Assumption}
\newtheorem{pro}{問題}

\newtheorem*{theorem*}{Theorem}
\newtheorem*{lemma*}{Lemma}
\newtheorem*{corollary*}{Corollary}
\newtheorem*{prop*}{Proposition}
\newtheorem*{rem*}{Remark}
\newtheorem*{assumption*}{Assumption}
\newtheorem*{notation*}{Notation}

\DeclareMathOperator{\spt}{spt} 

\newcommand{\mres}{\mathbin{\vrule height 1.6ex depth 0pt width
0.13ex\vrule height 0.13ex depth 0pt width 1.3ex}}
\renewcommand{\labelenumi}{(\arabic{enumi})}

\makeatletter
\c@MaxMatrixCols=11
\makeatother

\numberwithin{equation}{section}

\setlength{\parindent}{0pt}

\title{第2回 補足問題}
\author{}
\date{}

\begin{document}

\maketitle

\vspace{-1.5cm}

\prob 関数$f : (0, 1)\rightarrow \mathbb{R}$は $\dsp\lim_{x\rightarrow 0}f(x)=0$かつ$\dsp\lim_{x\rightarrow 0}\frac{f(x)-f(x/2)}{x}=0$を満たすとする. このとき, 極限値$\dsp\lim_{x\rightarrow 0}\frac{f(x)}{x}$を求めよ.

\prob $I$は区間, $f:I\rightarrow\mathbb{R}$は単調とし, 任意の$a, b\in I$と, $f(a)$と$f(b)$の間の任意の実数$\xi$に対し, ある$x\in I$が存在して $f(x)=\xi$が成立するとする. このとき$f$は連続であることを示せ. また, この問題は何が言いたいと思う？

\prob\label{aa} 連続関数 $f : \mathbb{R}\rightarrow \mathbb{R}$は任意の$x, y\in\mathbb{R}$に対して$|x-y|\leqslant|f(x)-f(y)|$を満たすとする.

(1) $f$は狭義単調であることを示せ.

(2) $f$は全射であることを示せ.

\prob 関数$f:[0, \infty)\rightarrow[0, \infty)$で, 任意の$x, y\in[0,\infty)$に対して
\[f(x)f(yf(x))=f(x+y)\]
を満たすものを全て答えよ.

\prob $f$ が $[0,\infty)$ 上で連続とする. ある $a>0$ が存在して $f$ が $[a,\infty)$ 上で一様連続であるとき, $f$ は $[0,\infty)$ 上で一様連続であることを示せ.

\prob $f$ を有界開区間 $(a , b)$ で連続な関数とする. 次の(1)と(2)が同値であることを示せ.

(1) $f$ は $(a , b)$ で一様連続である.

(2) $[a , b]$ 上の連続関数 $g$ で, すべての $x\in (a , b)$ に対して $g(x)=f(x)$ を満たすものが存在する.

\prob $\dsp f(x)=\frac{1}{x}$ は $(1 , \infty)$ でリプシッツ連続であるが, $(0 , \infty)$ でリプシッツ連続でないことを示せ.

\prob リプシッツ連続ならば一様連続であることを示せ. また一様連続であるがリプシッツ連続でない関数の例を挙げよ.

\prob $f(x)=
\begin{cases}
	\dsp\frac{1}{\log x} & (x\in (0,1/2]) \\ 0 & (x=0)
\end{cases}$ 
は一様連続であるが, 任意の $\alpha \in (0,1)$ に対し $\alpha$-H\"{o}lder連続でないことを示せ.


\end{document}